\begin{abstract}
    This study evaluates retrieval augmented generation (RAG) for medical question answering (MQA) by comparing a small language model (SLM), a large language model (LLM), and a rule based hybrid routing strategy using a 115 document biomedical corpus and open MQA benchmarks 9,090 predictions per configuration. Primary answer quality evidence comes from these open benchmarks MedQA, MMLU-Med, and PubMedQA, where RAG effects were task-dependent: MMLU Med performance improved for the LLM, while MedQA and PubMedQA label accuracy did not consistently increase with RAG. To assess deployment feasibility, we introduce an Energy Infrastructure Economy (EIE) framework that reports model quality alongside energy, carbon, and infrastructure costs. NVIDIA Management Library (NVML) based measurements on NVIDIA T4 GPUs indicate that SLM+RAG consumed approximately 2.15$\times$ less energy and produced 2.14$\times$ lower CO$_2$e emissions per million queries compared to LLM+RAG 19.82 vs.\ 42.47\,kg CO$_2$e. Automated faithfulness scores $\approx$47--48\% indicated only partial evidence support, highlighting the need for retrieval refinement and claim verifier benchmarking. Within the evaluated corpus, models, and benchmark scope, these findings suggest that retrieval augmented SLMs can deliver comparable automated faithfulness under the tested configuration, with substantially lower deployment cost and environmental impact than LLM+RAG. Clinical validation by practicing physicians is planned as follow up work; this manuscript does not present real world clinical safety or effectiveness results.
\end{abstract}
